\documentclass[11pt]{article}
\usepackage[margin=1in]{geometry}
\usepackage{amsmath,amssymb,amsthm,hyperref}
\hypersetup{colorlinks=true,urlcolor=blue,linkcolor=blue}
\newtheorem{theorem}{Theorem}
\newtheorem{lemma}{Lemma}
\newtheorem{corollary}{Corollary}
\newcommand{\I}{\mathcal I}
\newcommand{\T}{\mathcal T}
\newcommand{\A}{\mathcal A}
\newcommand{\B}{\mathcal B}
\title{A strictly positive biquadratic $M$-tensor\protect\\without a sum-of-squares representation\protect\\\large MathDB 371784}
\author{Principia Math In Progress}
\date{21 September 2026}
\begin{document}
\maketitle

\section*{Task 1: full resolution of the stated conjecture}
Xu, Cui and Qi ask whether every symmetric biquadratic $M$-tensor
is a sum of squares of bilinear forms \cite[Question 1]{XCQ}.
This was Conjecture 3.1 in their current preprint.
We answer negatively, even with a strict $M$-eigenvalue margin.

Here symmetric means the tensor identities
\begin{equation}\label{eq:sym}
                 a_{ijkl}=a_{kjil}=a_{klij}.
\end{equation}
It does not impose invariance of the polynomial under arbitrary
permutations of the $x$-variables or the $y$-variables. The source
explicitly distinguishes these notions. We use exactly
\eqref{eq:sym}, its hypothesis in Question 1.

For $x,y\in\mathbb R^3$ define
\begin{align}
 C(x,y)&=\sum_{i=1}^3x_i^2y_i^2+
             x_1^2y_2^2+x_2^2y_3^2+x_3^2y_1^2
             -2\sum_{i<j}x_ix_jy_iy_j,\label{eq:choi}\\
 J(x,y)&=\Bigl(\sum_{i=1}^3x_i^2\Bigr)
              \Bigl(\sum_{j=1}^3y_j^2\Bigr),\qquad
 P(x,y)=8C(x,y)+J(x,y).\label{eq:P}
\end{align}
The form $C$ is the classical Choi form \cite{Choi}, up to cyclic
orientation. We do not claim its PSD/non-SOS phenomenon as new.
The point here is its application to the stated $M$-tensor question,
with an explicit strict perturbation and a self-contained certificate.

\begin{theorem}\label{thm:main}
The symmetric biquadratic tensor $\A$ of $P$ satisfies
\[
 \A=9\I-\B,\qquad \B\ge0\text{ entrywise},\qquad
                   \lambda_{\max}(\B)=8<9.
\]
Thus $\A$ is a strong $M$-biquadratic tensor. Moreover,
$P(x,y)\ge\|x\|^2\|y\|^2$, but $P$ is not a sum of squares
of bilinear forms.
\end{theorem}

\subsection*{Positivity and the tensor entries}
Let $\T$ be the tensor of $C$. Its nonzero entries are as follows.
For $i=1,2,3$, with $i+1$ taken cyclically,
\[
 t_{iiii}=t_{i,i+1,i,i+1}=1.
\]
For each $i<j$ set
\[
 t_{iijj}=t_{jiij}=t_{ijji}=t_{jjii}=-\tfrac12.
\]
Every other entry is zero. These prescriptions verify
\eqref{eq:sym} and give
$C=\sum_{i,j,k,l}t_{ijkl}x_iy_jx_ky_l$.
The $M$-identity is $I_{ijkl}=\delta_{ik}\delta_{jl}$,
whose form is $J$. Consequently $\A=8\T+\I$ has integer entries.

\begin{lemma}\label{lem:psd}
$C(x,y)\ge0$ for all real $x,y$.
\end{lemma}
\begin{proof}
Write $C=x^TM(y)x$, where
\[
 M(y)=\begin{pmatrix}
 y_1^2+y_2^2&-y_1y_2&-y_1y_3\\
 -y_1y_2&y_2^2+y_3^2&-y_2y_3\\
 -y_1y_3&-y_2y_3&y_3^2+y_1^2
 \end{pmatrix}.
\]
Put $a=y_1^2,b=y_2^2,c=y_3^2$. Its diagonal entries are
nonnegative. The principal minors on index sets $12,13,23$ are,
respectively,
\[
 ac+b^2+bc,\qquad a^2+ab+bc,\qquad ab+ac+c^2.
\]
Its determinant is
\[
                a^2c+ab^2+bc^2-3abc\ge0,
\]
by AM--GM, since the product of the three positive terms is
$(abc)^3$. The inequality includes zero values by continuity.
All principal minors of the real symmetric matrix $M(y)$ are
nonnegative, so it is positive semidefinite. This proves the lemma.
\end{proof}

Set $\B=8(\I-\T)$. Diagonal entries of $\T$ are zero or one,
and its off-diagonal entries are zero or $-1/2$.
Therefore $\B$ has entries in $\{0,4,8\}$ and is nonnegative.
Its form is $8(J-C)$. If $\lambda$ is any $M$-eigenvalue of
$\B$ with unit eigenvectors $x,y$, contraction with $x$ gives
\[
              \lambda=8(1-C(x,y))\le8.
\]
To verify attainment directly, set $v=(1,1,1)^T/\sqrt3$.
For each fixed first index, the entries of $\B$ sum to $24$;
the same holds for each fixed second index. Hence
\[
       \B\mathbin{\cdot}vvv=8v,\qquad
       \B v\mathbin{\cdot}vv=8v.
\]
Thus $8$ is an $M$-eigenvalue and $\lambda_{\max}(\B)=8$.
This proves the asserted strict $M$-tensor condition.
Lemma~\ref{lem:psd} also gives $P\ge J$.

\subsection*{An exact obstruction to every SOS representation}
Define a linear functional $\ell$ on the space of real biquadratic
polynomials by specifying its values on monomials:
\[
\begin{array}{c|c}
\text{monomial}&\ell(\text{monomial})\\\hline
x_i^2y_i^2&2\\
x_1^2y_2^2,\ x_2^2y_3^2,\ x_3^2y_1^2&1\\
x_2^2y_1^2,\ x_3^2y_2^2,\ x_1^2y_3^2&4\\
x_ix_jy_iy_j\quad(i<j)&2
\end{array}
\]
Its value on every unlisted monomial is zero.
For an arbitrary bilinear form $g=\sum_{i,j=1}^3u_{ij}x_iy_j$,
direct expansion gives
\begin{align}
 \ell(g^2)={}&2(u_{11}+u_{22}+u_{33})^2
        +(u_{12}+2u_{21})^2\notag\\
       &+(u_{23}+2u_{32})^2+(u_{31}+2u_{13})^2\ge0.
       \label{eq:separator}
\end{align}
In particular $\ell$ is nonnegative on every finite sum of squares
of bilinear forms. On the other hand,
\[
 \ell(C)=6+3-12=-3,\qquad \ell(J)=6+3+12=21,
 \qquad \ell(P)=8(-3)+21=-3<0.
\]
This contradicts any proposed SOS representation of $P$ and completes
the proof of Theorem~\ref{thm:main}.

\begin{corollary}
For every $0<t<1/7$, the form $C+tJ$ is strictly positive on
pairs of nonzero vectors, its tensor is a strong $M$-biquadratic
tensor, and it is not SOS.
\end{corollary}
\begin{proof}
Its tensor is $(1+t)\I-(\I-\T)$, with
$\lambda_{\max}(\I-\T)=1$. Its form is at least $tJ$, while
$\ell(C+tJ)=-3+21t<0$.
\end{proof}

\paragraph{Verification and scope.}
The accompanying verifier checks all tensor entries and symmetries,
all principal-minor identities, the exact eigenvector equations,
and the separator identity for nine symbolic bilinear coefficients.
These checks support the displayed algebraic proof; no numerical
optimization, eigenvalue approximation, or sampled positivity claim
is used. This resolves Question 1 only. Review is by the authoring
agent, without independent peer review or a proof assistant.

\section*{Four proposed follow-up tasks}
\begin{enumerate}
\item Task 2: Determine the exact smallest $t\ge0$ for which $C+tJ$
is SOS, and exhibit a representation at the threshold.
\item Task 3: Classify the SOS members of the three-parameter cyclic
family with independently weighted diagonal and two cyclic square sums.
\item Task 4: Characterize which symmetric $M$-biquadratic tensors
admit positive diagonal rescaling to a diagonally dominated tensor.
\item Task 5: Obtain dimension-dependent bounds for the amount of
$M$-identity that must be added to a normalized $M$-biquadratic tensor
to make it SOS.
\end{enumerate}
These are proposed research tasks and earn no additional solution credit.

\begin{thebibliography}{9}
\bibitem{XCQ} Y.~Xu, C.~Cui and L.~Qi,
\emph{Sum of Squares Decompositions for Structured Biquadratic Forms},
Communications on Applied Mathematics and Computation (2026),
\href{https://doi.org/10.1007/s42967-026-00586-7}{doi:10.1007/s42967-026-00586-7}.
Preprint: \href{https://arxiv.org/abs/2512.02734v2}{arXiv:2512.02734v2}.
\bibitem{Choi} M.-D.~Choi,
\emph{Positive semidefinite biquadratic forms},
Linear Algebra and its Applications \textbf{12} (1975), 95--100,
\href{https://doi.org/10.1016/0024-3795(75)90058-0}{doi:10.1016/0024-3795(75)90058-0}.
\end{thebibliography}
\end{document}
